% ============================================================
%  03_state_of_the_art.tex  –  Chapter 2: State of the Art
% ============================================================

\chapter{State of the Art}
\label{ch:sota}

\section{Air Quality Monitoring Systems}
\label{sec:aq-monitoring}

Indoor air quality monitoring has become an increasingly important research topic, driven by
growing evidence linking extended exposure to indoor pollutants with chronic health conditions
such as respiratory infections, cardiovascular disease, and lung cancer. IoT technologies are
now used to provide real-time data collection, at a more affordable price.

The parameters which are most monitored are temperature, humidity, CO\textsubscript{2}, carbon
monoxide, particulate matter and volatile organic compounds. Used sensors are usually low-cost,
connected to Arduino or Raspberry Pi units, which have the role of processing units. A review
by Marques (2020) of 40 IoT-based indoor air quality systems showed that 35\% used
Raspberry Pi, confirming it as one of the most popular platforms for this type of
application~\cite{marques2020}. From the same review we find out that 72.5\% of systems claim
energy efficiency as a design goal. Also the fact that cloud connection is favored for data
storage and remote access.

The main issue in this field is finding the balance between sensor coverage, costs and data
reliability. Ganesan et al.\ (2021) demonstrated that a Raspberry Pi-based system monitoring
10 simultaneous environmental parameters, including CO, CO\textsubscript{2}, NO\textsubscript{2},
SO\textsubscript{2}, O\textsubscript{3}, and TVOCs, can be built at low cost and deployed in
real residential and office environments~\cite{ganesan2021}. This showed that monitoring
multiple pollutants was not only possible from a consumer standpoint, but also that different
environments have different levels of pollutants, resulting in the importance of local
monitoring. Likewise, Kim et al.\ (2024) successfully developed and tested a low-cost system
using Raspberry Pi to measure 11 environmental parameters and succeeded in obtaining stable
measurements from sensors after testing for statistics at a cost of around 94 USD worth of
hardware~\cite{kim2024}.

A popular trend in this field is now the use of prediction capabilities in monitoring systems.
Fahmy et al.\ (2021) suggested a hybrid deep learning system using CNN and LSTM layers to
predict hourly air quality with the optimization performed for Raspberry Pi
4~\cite{fahmy2021}. Later, Gardner et al.\ (2024) looked into TinyML approaches to develop
light-weight prediction models to run on microcontrollers that are capable of predicting air
quality despite having limited resources~\cite{gardner2024}. The mentioned publications provided
inspiration for the prediction component of this project.

\section{Similar Works}
\label{sec:similar-works}

Currently, several commercial and research-based air quality monitoring systems are available.
These range from products destined for consumers to research-oriented IoT platforms. Each type
having their own features and limitations.

The \textbf{Awair Element} is a product type indoor air quality monitor. It tracks values of
temperature, humidity, CO\textsubscript{2}, TVOCs and PM2.5. It connects to a smartphone
application and displays real-time readings and an air quality score. It doesn't offer the
option to export the raw data read or to make changes to the sensors~\cite{awair}.

\textbf{uRADMonitor A3} is a more advanced product device. Its target audience are consumers
and researchers. It monitors a wider range of pollutants and it sends data to a central online
platform. It supports data export, but it only offers limited hardware modification possibility
or local data hosting~\cite{uradmonitor}.

Below are characteristics specific to these two products and also characteristics of this
project, all side by side:

\vspace{12pt}

\begin{table}[H]
  \centering
  \caption{Table 1. System characteristics comparison}
  \label{tab:comparison}
  \renewcommand{\arraystretch}{1.3}
  \begin{tabularx}{\textwidth}{|l|X|X|X|}
    \hline
    \textbf{Feature} & \textbf{Awair Element} & \textbf{uRADMonitor A3} & \textbf{Current Project} \\
    \hline
    Hardware platform
      & Proprietary
      & Proprietary
      & Raspberry Pi 4 \\
    \hline
    Monitored parameters
      & Temp, humidity, CO\textsubscript{2}, TVOC, PM2.5
      & PM, O\textsubscript{3}, HCHO, and others
      & Temp, humidity, CO, LPG, alcohol, H\textsubscript{2}, air quality, sound \\
    \hline
    Web interface    & No  & Online platform & Yes \\
    \hline
    Mobile interface & Yes & No  & Yes \\
    \hline
    Data export (CSV)& No  & Limited & Yes \\
    \hline
    ML prediction    & No  & No  & Yes \\
    \hline
    Open / customizable & No & Partial & Yes \\
    \hline
    Cost & \textasciitilde\$299 & \textasciitilde\$350 & Below \$100 \\
    \hline
  \end{tabularx}
\end{table}
